\documentclass[10pt]{article}

% ---------- Document Identity (update these only) ----------
\newcommand{\DocStage}{}
\newcommand{\DocTitle}{Your Road to Tier One SOF}
\newcommand{\ProgramName}{182-Week Civilian-to-Selection Readiness Pipeline}
\newcommand{\ProgramSubtitle}{}

% ---------- Page ----------
\usepackage[legalpaper,left=0.5in,right=0.5in,top=0.6in,bottom=0.45in]{geometry}

% ---------- Typography ----------
\usepackage[T1]{fontenc}
\usepackage[utf8]{inputenc}
\usepackage{libertinus}
\usepackage{microtype}
\usepackage{parskip}
\usepackage{ragged2e}
\usepackage{textcomp}
\AtBeginDocument{\color{black}}
\AtBeginDocument{\pagecolor{white}}
\AtBeginDocument{\justifying}

% ---------- Landscape pages ----------
\usepackage{pdflscape}

% ---------- Color ----------
\usepackage[table]{xcolor}
\definecolor{StageBlue}{HTML}{1F4E79}
\definecolor{StageBlueDark}{HTML}{163A5A}
\definecolor{LightGray}{HTML}{F2F4F7}
\definecolor{MidGray}{HTML}{D9DEE5}
\definecolor{RuleGray}{HTML}{C7CED8}
\definecolor{HeaderInk}{HTML}{000000}
\definecolor{Ink}{HTML}{000000}

% ---------- Utility ----------
\usepackage{enumitem}
\sloppy
\setlength{\emergencystretch}{2em}

% ---------- Boxes / UI ----------
\usepackage[most]{tcolorbox}

\tcbset{
  charterTitle/.style={
    enhanced,
    colback=StageBlue!4,
    colframe=StageBlue,
    boxrule=1.25pt,
    arc=3pt,
    left=16pt,right=16pt,top=14pt,bottom=14pt,
    borderline north={3.2pt}{0pt}{StageBlue},
    borderline west={4pt}{0pt}{StageBlue},
    borderline south={1.25pt}{0pt}{StageBlue},
    drop shadow={black!25!white},
  },
  charterRule/.style={
    enhanced,
    colback=LightGray,
    colframe=RuleGray,
    boxrule=0.85pt,
    arc=3pt,
    left=12pt,right=12pt,top=10pt,bottom=10pt,
    borderline west={3pt}{0pt}{StageBlue},
  },
  charterBadge/.style={
    enhanced,
    colback=StageBlue,
    coltext=white,
    colframe=StageBlueDark,
    boxrule=0.6pt,
    arc=2pt,
    left=6pt,right=6pt,top=3pt,bottom=3pt,
  }
}
\newtcbox{\Badge}{nobeforeafter,charterBadge}

% ---------- Lists ----------
\setlist[itemize]{leftmargin=1.5em, itemsep=0.25em, topsep=0.25em}
\setlist[enumerate]{leftmargin=1.8em, itemsep=0.25em, topsep=0.25em}

% ---------- TOC ----------
\usepackage{tocloft}
\setcounter{tocdepth}{3}
\setlength{\cftbeforesecskip}{3pt}
\setlength{\cftbeforesubsecskip}{2pt}
\setlength{\cftbeforesubsubsecskip}{0pt}
\renewcommand{\cftsecfont}{\bfseries}
\renewcommand{\cftsecpagefont}{\bfseries}
\renewcommand{\cftsubsecfont}{\normalfont}
\renewcommand{\cftsubsecpagefont}{\normalfont}
\renewcommand{\cftsubsubsecfont}{\normalfont}
\renewcommand{\cftsubsubsecpagefont}{\normalfont}
\setlength{\cftsecindent}{0pt}
\setlength{\cftsubsecindent}{1.25em}
\setlength{\cftsubsubsecindent}{2.8em}
\setlength{\cftsecnumwidth}{2.1em}
\setlength{\cftsubsecnumwidth}{2.8em}
\setlength{\cftsubsubsecnumwidth}{3.6em}

% Remove the built-in TOC heading so only the custom heading is shown
\makeatletter
\renewcommand{\tableofcontents}{\@starttoc{toc}}
\makeatother

% ---------- Header/Footer: page number only ----------
\usepackage{fancyhdr}
\pagestyle{fancy}
\fancyhf{}
\renewcommand{\headrulewidth}{0pt}
\renewcommand{\footrulewidth}{0pt}
\fancyfoot[C]{\thepage}

% ---------- Section formatting ----------
\usepackage{titlesec}

\titleformat{\section}
  {\Large\bfseries\color{Ink}}
  {\thesection}{0.6em}{}
  [\vspace{0.25em}\textcolor{StageBlue}{\rule{\linewidth}{0.8pt}}]

\titleformat{\subsection}
  {\large\bfseries\color{Ink}}
  {\thesubsection}{0.6em}{}

\titleformat{\subsubsection}
  {\normalsize\bfseries\color{Ink}}
  {\thesubsubsection}{0.6em}{}

\titlespacing*{\section}{0pt}{0.95em}{0.35em}
\titlespacing*{\subsection}{0pt}{0.65em}{0.25em}
\titlespacing*{\subsubsection}{0pt}{0.55em}{0.2em}

% ---------- Table engine ----------
\usepackage{tabularray}
\UseTblrLibrary{booktabs}

% ---------- tabularray: GLOBAL table treatment ----------
\SetTblrInner{
  cells = {fg=black, bg=white, valign=t},
}

% ---------- Header helpers ----------
\newcommand{\Hdr}[1]{\shortstack[c]{\strut #1\strut}}

% ---------- Lines ----------
\newcommand{\TitleLine}{\textcolor{StageBlue}{\rule{0.98\linewidth}{0.95pt}}}
\newcommand{\SubTitleLine}{\textcolor{RuleGray}{\rule{0.98\linewidth}{0.65pt}}}

% ---------- Continued on next page ----------
\newcommand{\ContinuedOnNextPageInline}{%
  \par
  \vfill
  \noindent\textcolor{RuleGray}{\rule{\linewidth}{1pt}}\vspace{-2pt}
  \noindent\hfill\textit{Continued on next page}\par\vspace{-10pt}
  \noindent\textcolor{RuleGray}{\rule{\linewidth}{1pt}}
  \clearpage
}

% ---------- PDF hyperlinks ----------
\usepackage[hidelinks]{hyperref}
\hypersetup{
  colorlinks=false,
  pdfborder={0 0 0},
  linktoc=all,
  pdftitle={\DocTitle},
  pdfauthor={\ProgramName}
}

% =========================================================
\begin{document}

\section*{7.15 — General Training and Daily Wear Apparel}
\phantomsection
\addcontentsline{toc}{section}{7.15 — General Training and Daily Wear Apparel}

\subsection*{7.15.1 Purpose \& Scope}
\phantomsection
\addcontentsline{toc}{subsection}{7.15.1 Purpose \& Scope}

This overview covers both:

\begin{itemize}
\item Training apparel for routine sessions (comfort, mobility, moisture, friction control).
\item Daily-life apparel support for start-from-zero adherence (baseline clothing posture that reduces friction and supports consistent execution).
\end{itemize}

This overview crosswalks Stage 6 timing to the apparel conditions required for safe routine execution and admissible evidence.

This category does not cover, except by cross-reference:

\begin{itemize}
\item Stage 7.14 environmental / extreme-weather systems,
\item Stage 7.11 footwear / footcare,
\item Stage 7.19 protective / visibility gear,
\item Stage 7.20 hygiene / consumables,
\end{itemize}

This overview does not provide programming, test protocols, calendars, medical guidance, or consumable anti-chafe prescriptions.

\subsection*{7.15.2 Governance And Safety Boundaries}
\phantomsection
\addcontentsline{toc}{subsection}{7.15.2 Governance And Safety Boundaries}

\begin{itemize}
\item Stage 0 boundary alignment: stop rules, safety override doctrine, conflict-resolution hierarchy, and evidence-admissibility controls apply to all 7.15 selections and substitutions.
\item Clothing affects compliance and injury risk through chafe, moisture retention, and overheating posture.
\item Start-from-zero posture: minimum viable clothing exists early without excessive spend; scale rotation depth as volume increases.
\item Validity posture:
  \begin{itemize}
  \item stable apparel reduces confounding; avoid novelty near protected windows.
  \item document major apparel system changes near gate windows at a high level (e.g., “new fabric system introduced”) if unavoidable.
  \end{itemize}
\item Evidence inadmissibility: evidence produced while apparel instability, unsafe chafe escalation, or unmanaged wet-clothing friction materially affects execution is inadmissible until corrected and re-verified.
\item Anti-chafe posture is governance-only here:
  \begin{itemize}
  \item adjust fit/fabric first,
  \item consumables live in Stage 7.20 (cross-reference only).
  \end{itemize}
\end{itemize}

\subsection*{7.15.3 Tier Doctrine}
\phantomsection
\addcontentsline{toc}{subsection}{7.15.3 Tier Doctrine}

\begin{itemize}
\item Price band convention: \$ = minimal household-cost item; \$\$ = modest purpose-built item; \$\$\$ = expanded-capability or redundancy item.
\item N/A rule: use N/A only when a field is genuinely not applicable and omission does not obscure safety, maintenance, or phase-timing logic.
\item Substitution posture: substitutions are acceptable only when they preserve fit, friction control, moisture management, and adherence reliability; substitutes that increase chafe, overheating, or repeated wet re-use break intent.
\item Tier 0: constraint state; household-only clothing posture. Minimal and not reliable for repeated training. Used only as a temporary bridge.
\item Tier 1: minimum viable kit with enough rotation depth to support frequent training and basic daily life without constant laundry crises; includes go-bag clothing kit.
\item Tier 2: expanded rotation that reduces laundry pressure and provides fast-dry spares as sweat volume increases.
\item Tier 3: staged kits (home/vehicle/travel) and redundancy to minimize variance and prevent missed sessions in high-consequence phases.
\end{itemize}

\subsection*{7.15.4 Selection Criteria \& Fit Checks}
\phantomsection
\addcontentsline{toc}{subsection}{7.15.4 Selection Criteria \& Fit Checks}

Fit and safety checklist (operational):

Fit rules:

\begin{itemize}
\item Full range of motion without binding: squat/hinge/arm swing should not be restricted.
\item Seam pressure points: identify common rub points (inner thigh, underarm, waistband) and remove/replace items that create repeatable hotspots.
\end{itemize}

Fabric rules:

\begin{itemize}
\item Moisture and drying posture: fast-dry fabrics reduce chafe and improve comfort in high-volume phases.
\item Cotton caution: cotton retains moisture and increases chafe risk during long, sweaty sessions; not prohibited, but treated as higher friction risk.
\end{itemize}

Chafe-control posture (non-medical):

\begin{itemize}
\item Adjust fit and fabric first (short length, liner fit, seam placement).
\item If chafe persists, cross-reference Stage 7.20 for consumable anti-chafe options; do not list consumables here.
\end{itemize}

Laundry feasibility:

\begin{itemize}
\item Rotation counts must be tied to realistic wash schedule and drying time. If laundry frequency cannot support six sessions per week, increase rotation depth rather than skipping sessions.
\end{itemize}

Rotation Targets by Tier Mini-Block (training and daily wear)

\begin{itemize}
\item Tier 1 targets:
  \begin{itemize}
  \item Training: 3–5 shirts; 2–4 bottoms; 7–10 underwear; 7–10 socks.
  \item Daily wear: 3–5 outfits.
  \end{itemize}
\item Tier 2 targets:
  \begin{itemize}
  \item Training: 6–8 shirts; 3–5 bottoms; 10–14 underwear; 10–14 socks; 2–3 fast-dry spares.
  \item Daily wear: 6–8 outfits.
  \end{itemize}
\item Tier 3 targets:
  \begin{itemize}
  \item Training: 10+ shirts; 5+ bottoms; 14+ underwear; 14+ socks; staged duplicates.
  \item Daily wear: 8+ outfits plus staged kits.
  \end{itemize}
\end{itemize}

\subsection*{7.15.5 Setup, Safety, and Anchoring}
\phantomsection
\addcontentsline{toc}{subsection}{7.15.5 Setup, Safety, and Anchoring}

Setup doctrine (no programming):

\begin{itemize}
\item Pre-session checks:
  \begin{itemize}
  \item remove snag hazards, abrasive zippers, and pocket items that rub.
  \item verify drawstrings and waistbands are secure without restricting breathing.
  \end{itemize}
\item Chafe prevention workflow posture:
  \begin{itemize}
  \item select known “no-hotspot” kit for longer sessions,
  \item change out of wet clothes promptly after sessions to reduce skin irritation.
  \end{itemize}
\item Cross-reference Stage 7.19 for visibility / traffic gear; do not rely on Stage 7.15 for visibility.
\end{itemize}

\subsection*{7.15.6 Maintenance, Replacement, and Storage}
\phantomsection
\addcontentsline{toc}{subsection}{7.15.6 Maintenance, Replacement, and Storage}

\begin{itemize}
\item Laundry cadence posture:
  \begin{itemize}
  \item maintain enough rotation to avoid repeated damp re-use.
  \item staged clean kit posture reduces missed sessions.
  \end{itemize}
\item Replacement cues:
  \begin{itemize}
  \item seam failure, persistent odor that cannot be removed, loss of elasticity, repeated chafe hotspots.
  \end{itemize}
\item Storage workflow:
  \begin{itemize}
  \item staged clean training kit and daily kit.
  \item wet / dirty separation posture via cross-reference Stage 7.10 / Stage 7.20.
  \end{itemize}
\end{itemize}

\subsection*{7.15.7 Substitution Rules — Minimal-Path Alternatives}
\phantomsection
\addcontentsline{toc}{subsection}{7.15.7 Substitution Rules — Minimal-Path Alternatives}

\begin{itemize}
\item If clothing is limited:
  \begin{itemize}
  \item prioritize clean/dry base layers and socks before any longer exposure.
  \item adjust modality to reduce friction when feasible (e.g., indoor low-impact work rather than long wet outdoor exposure).
  \end{itemize}
\item Do not use clothing that increases chafe or overheating on longer sessions; substitute to a safer kit rather than pushing through.
\item Document major apparel changes near tests (high level) to preserve comparability: “new fabric system” or “new fit system” as a controlled variable.
\end{itemize}

\subsection*{7.15.8 Phase Integration Map — Required}
\phantomsection
\addcontentsline{toc}{subsection}{7.15.8 Phase Integration Map — Required}

\begingroup
\scriptsize
\setlength{\tabcolsep}{2pt}
\renewcommand{\arraystretch}{1.10}

\begin{longtblr}{
  width=\linewidth,
  rowhead=1,
  colspec = {
    Q[l,wd=1.05in]
    Q[l,wd=1.55in]
    X[l,2.75]
    X[l,3.65]
  },
  hlines,
  vlines,
  cells = {hsep=2pt, vsep=6pt, halign=l, valign=t},
  row{1} = {bg=StageBlue!15, fg=black, font=\bfseries, halign=l, valign=m},
  row{odd} = {bg=white},
  row{even} = {bg=LightGray},
}
\Hdr{Phase} &
\Hdr{Required Tier (from Stage 6)} &
\Hdr{What Changes / Becomes Mandatory} &
\Hdr{Notes} \\
Phase 00 &
Tier 1 &
Tier 1 minimum viable training and daily-wear rotation becomes mandatory; go-bag kit established by Phase 00 (End) &
Phase 00 (End): minimum rotation prevents immediate compliance collapse. Phase 00 (End): go-bag kit supports continuity; cross-reference Stage 7.10 staging. \\
Phase 01 &
Tier 1 &
Maintain Tier 1; stabilize kit choices for protected windows &
Phase 01 (End): avoid novelty near Deload+Test windows; manage sweat/chafe. \\
Phase 02 &
Tier 1 &
Maintain Tier 1; friction control becomes more consequential as workload rises &
Phase 02 (End): ensure enough rotation to avoid re-wearing damp kit. \\
Phase 03 &
Tier 1 &
Maintain Tier 1; ruck exposure begins to change friction points &
Phase 03 (End): ruck straps increase rub points; adjust fit/fabric first; cross-reference Stage 7.11 and Stage 7.20 for consumables. \\
Phase 04 &
Tier 1 &
Maintain Tier 1; wet exposure increases gear turnover &
Phase 04 (End): fast drying becomes more valuable; consider Tier 2 if laundry pressure is failing. \\
Phase 05 &
Tier 1 &
Maintain Tier 1; longer durations increase sweat load &
Phase 05 (End): avoid cotton-heavy kit for longer sweaty sessions if chafe/irritation occurs. \\
Phase 06 &
Tier 1 &
Tier 2 is not required by Stage 6 for 7.15; maintain Tier 1 minimum; optionally expand rotation if laundry pressure fails &
Phase 06 (End): Stage 6 keeps Stage 7.15 at Tier 1; if high volume makes Tier 1 insufficient in practice, treat as execution risk and consider Stage 8 review only if it becomes a requirement. \\
Phase 07 &
Tier 1 &
Maintain Tier 1; consider expanded rotation to reduce variance &
Phase 07 (End): hybrid density increases sweat/friction; operationally, Tier 2 may reduce misses but remains optional unless Stage 6 changes. \\
Phase 08 &
Tier 1 &
Maintain Tier 1; load tolerance phases increase wet/dirty turnover &
Phase 08 (End): prioritize staging and drying via Stage 7.10; hygiene via Stage 7.20. \\
Phase 09 &
Tier 1 &
Maintain Tier 1; high \textbf{CAL} volume increases sweat and friction &
Phase 09 (End): expanded rotation is high ROI but not required unless Stage 6 changes. \\
Phase 10 &
Tier 1 &
Maintain Tier 1; multi-domain gates increase consequence of skin breakdown &
Phase 10 (End): select stable “no-hotspot” kit for decision windows; no novelty. \\
Phase 11 &
Tier 1 &
Maintain Tier 1; recovery sensitivity increases consequence of irritation &
Phase 11 (End): avoid preventable skin breakdown; keep kit clean-to-ready. \\
Phase 12 &
Tier 1 &
Maintain Tier 1; Stage 6 does not elevate 7.15 above Tier 1 &
Phase 12 (End): if staged kits (Tier 3) are adopted for reliability, they remain optional unless Stage 6 changes; avoid introducing novel fabrics in protected windows. \\
Phase 13 &
Tier 1 &
Maintain Tier 1; consolidation requires low variance &
Phase 13 (End): maintain stable kit; no novelty near final protected windows. \\
\end{longtblr}

\endgroup

7.15.8 Phase Integration Map Notes:

\begin{itemize}
\item Stage 6 baseline keeps Stage 7.15 at Tier 1 across all phases. If any phase artifact requires Tier 2+ apparel rotation as mandatory, requires Stage 8 review.
\end{itemize}

\subsection*{7.15.9 Risks, Contraindications, and Red Flags}
\phantomsection
\addcontentsline{toc}{subsection}{7.15.9 Risks, Contraindications, and Red Flags}

Non-medical risks:

\begin{itemize}
\item Chafing leading to altered mechanics (gait changes, reduced ROM, avoidance behaviors).
\item Overheating from poor fabric choice and lack of venting; increases heat stress risk.
\item Skin irritation from detergents or products (cross-reference Stage 7.20); repeated irritation requires process changes and evaluation if severe or persistent.
\end{itemize}

Red flags (stop and seek evaluation posture):

\begin{itemize}
\item Severe skin breakdown, spreading redness or heat, drainage, or fever/systemic illness signs → stop exposure and seek evaluation posture.
\item Persistent rash or swelling that worsens with training → stop the provoking exposure and seek evaluation posture.
\end{itemize}

Requires Stage 8 review triggers:

\begin{itemize}
\item Any mismatch where a phase requires 7.15 higher than Tier 1 earlier than Stage 6.
\item Any attempt to include anti-chafe consumables in tier tables (must be Stage 7.20 cross-reference).
\item Any attempt to include gloves in tier tables (must cross-reference Stage 7.19 / Stage 7.01).
\end{itemize}

\subsection*{7.15.10 Compliance Checklist — Pass/Fail}
\phantomsection
\addcontentsline{toc}{subsection}{7.15.10 Compliance Checklist — Pass/Fail}

\begin{itemize}
\item \textbf{PASS}/\textbf{FAIL}: Tier 0–Tier 3 tables use the required schema, generic item types only, and include numeric count ranges for training and daily wear.
\item \textbf{PASS}/\textbf{FAIL}: Tier 1 rotation exists by Phase 00 (End) per Stage 6 and includes go-bag clothing kit (cross-referenced to Stage 7.10 staging).
\item \textbf{PASS}/\textbf{FAIL}: Anti-chafe consumables are not listed in tier tables; cross-reference to Stage 7.20 only.
\item \textbf{PASS}/\textbf{FAIL}: Laundry tools are not listed; cross-references to Stage 7.10 / Stage 7.20 are used.
\item \textbf{PASS}/\textbf{FAIL}: Gloves are not listed; cross-references to Stage 7.19 / Stage 7.01 are used.
\item \textbf{PASS}/\textbf{FAIL}: Phase Integration Map is included and aligned to Stage 6 timing; any mismatch is flagged for Stage 8 review.
\end{itemize}

\end{document}